\section{System Context and Trust Placement}\label{sec:system-context}

\AFD{} is an endpoint-local attachment transformation and release gate. Its
placement follows from the confidentiality objective of end-to-end encrypted
messaging. On the sending endpoint, an attachment is reconstructed and
validated before the messaging application encrypts it. On the receiving
endpoint, ciphertext is authenticated and decrypted before the resulting
attachment bytes enter the same defense pipeline. The encrypted relay carries
the protocol object but does not receive attachment plaintext, choose a media
handler, or issue the final delivery verdict. This arrangement preserves the
blind-relay property expected of secure messaging systems~\cite{unger2015}.
It also avoids treating transport encryption as an answer to risks that arise
only when an endpoint parser consumes decrypted bytes.

The outgoing and incoming placements have complementary purposes. Outgoing
processing limits the representations that an honest sender injects into the
encrypted conversation. Incoming processing remains necessary because a
receiver cannot assume that every sender runs the same release, policy, or
implementation. It also covers imported ciphertext, compromised senders, and
protocol peers outside the local administrative domain. The two placements do
not require the relay to inspect content. They require each endpoint to bind
publication to its own authenticated policy snapshot and to run the gate on
the exact bytes that would otherwise reach storage, preview, or rendering.

\begin{figure}[t]
\centering
\resizebox{0.98\textwidth}{!}{%
\begin{tikzpicture}[
  node distance=7mm and 9mm,
  endpoint/.style={draw,rounded corners=2pt,minimum width=47mm,minimum height=31mm,fill=gray!6},
  proc/.style={draw,rounded corners=2pt,align=center,minimum width=28mm,minimum height=8mm,font=\small},
  relay/.style={draw,align=center,minimum width=32mm,minimum height=17mm,fill=blue!7,font=\small},
  arrow/.style={-{Latex[length=2mm]},thick},
  cipher/.style={-{Latex[length=2mm]},thick,dashed}
]
\node[endpoint] (sender) {};
\node[font=\small\bfseries,anchor=north] at ([yshift=-2mm]sender.north) {Sending endpoint};
\node[proc,anchor=south] (spre) at ([yshift=3mm]sender.south) {reconstruct, validate,\\clean-only release};
\node[proc,above=4mm of spre] (senc) {authenticated\\encryption};

\node[relay,right=18mm of sender] (relay) {blind relay\\ciphertext routing};

\node[endpoint,right=18mm of relay] (receiver) {};
\node[font=\small\bfseries,anchor=north] at ([yshift=-2mm]receiver.north) {Receiving endpoint};
\node[proc,anchor=south] (rpost) at ([yshift=3mm]receiver.south) {reconstruct, validate,\\clean-only publication};
\node[proc,above=4mm of rpost] (rdec) {authenticated\\decryption};

\draw[arrow] (spre) -- node[right,font=\scriptsize]{accepted bytes} (senc);
\draw[cipher] (senc.east) -- node[above,font=\scriptsize]{ciphertext} (relay.west);
\draw[cipher] (relay.east) -- node[above,font=\scriptsize]{ciphertext} (rdec.west);
\draw[arrow] (rdec) -- node[right,font=\scriptsize]{plaintext candidate} (rpost);
\end{tikzpicture}%
}
\caption{Endpoint placement around end-to-end encryption. The relay preserves
ciphertext routing and learns neither the reconstructed attachment nor the
local verdict. The receiving endpoint repeats the gate after decryption.}
\label{fig:system-endpoint-placement}
\end{figure}

The principal protected assets are endpoint execution integrity, client
availability, privacy-sensitive metadata, reconstructed-content integrity,
policy consistency, and attachment confidentiality in transit. Input bytes,
names, caller-supplied media types, dimensions, offsets, codec payloads, and
container metadata are adversary controlled. The authenticated application
binary, the selected local policy, the cryptographic transport implementation,
and the operating-system enforcement mechanisms are trusted. External
decoders and encoders are not treated as validators. Their output must pass a
separate canonical parser, although semantic reconstruction still depends on
their correct decoding and encoding behavior.

The architecture is intentionally conjunctive. The facade derives concrete
format evidence from the portable name, caller-supplied media type, and byte
recognizer. It checks compatibility before selecting a handler. The selected
policy fixes the minimum reconstruction level before dispatch. A handler then
returns candidate bytes, a declared output media type, the reconstruction
level actually performed, and stage evidence. The root pipeline independently
parses the candidate, checks output-type agreement and policy allowlists, scans
the rebuilt bytes for configured signatures, applies output-size and expansion
bounds, and derives a verdict. Only a clean verdict retains candidate bytes in
the public result.

\begin{figure}[t]
\centering
\resizebox{0.99\textwidth}{!}{%
\begin{tikzpicture}[
  node distance=6mm and 6mm,
  box/.style={draw,rounded corners=2pt,align=center,minimum width=24mm,minimum height=10mm,font=\small},
  control/.style={box,fill=blue!7},
  transform/.style={box,fill=green!8},
  gate/.style={box,fill=yellow!10},
  terminal/.style={box,fill=gray!12},
  arrow/.style={-{Latex[length=2mm]},thick}
]
\node[control] (policy) {immutable policy\\snapshot};
\node[control,right=of policy] (evidence) {name, media type,\\byte profile};
\node[transform,right=of evidence] (handler) {format handler\\candidate output};
\node[gate,right=of handler] (validate) {independent parse,\\type and limits};
\node[gate,right=of validate] (scan) {post-rebuild\\signature scan};
\node[terminal,right=of scan] (release) {clean bytes or\\empty output};
\draw[arrow] (policy) -- (evidence);
\draw[arrow] (evidence) -- (handler);
\draw[arrow] (handler) -- (validate);
\draw[arrow] (validate) -- (scan);
\draw[arrow] (scan) -- (release);
\end{tikzpicture}%
}
\caption{Reference component path. A handler creates a candidate but cannot
authorize delivery. Validation, signature policy, and final verdict derivation
remain in the facade.}
\label{fig:component-architecture}
\end{figure}

The current workspace divides these responsibilities across a root facade and
five supporting crates. \texttt{defender-core} contains the versioned policy
document and migration boundary. The media-handler support crate defines the
handler contract. The signature crate compiles baseline and
administrator rules into a multi-pattern scanner. \texttt{defender-ffi}
defines the versioned protobuf model used across the native boundary.
The observability crate defines audit-event and sink contracts. The root
\texttt{file\_defender} crate owns orchestration, concrete handlers,
canonical validation, tenant resolution, filesystem publication, and the C
ABI. This decomposition reduces accidental authority. In particular, neither
the signature crate nor a handler can write a deliverable file.

The threat vocabulary does not alter this boundary. Viruses, worms, Trojans,
ransomware, spyware, and downloaders identify threat and corpus strata. Their
carrier properties motivate measurements such as executable-header rejection,
metadata exclusion, polyglot closure, bounded parsing, and nonpublication of
unsupported objects. \AFD{} does not behaviorally classify those malware
families. It also does not prove that arbitrary accepted bytes are harmless.
An accepted artifact belongs to a constrained output grammar under the active
policy and has passed the implemented gates. Harmful visual meaning, steganic
information retained in semantic media, and vulnerabilities in downstream
renderers remain outside that statement.

Deployment therefore assumes that callers preserve gate ordering and honor
the returned deliverability condition. They must not encrypt an outgoing
candidate before a clean verdict, publish decrypted input while processing is
pending, or recover candidate bytes from diagnostics after a non-clean result.
The Rust, C, and Swift interfaces erase deliverable bytes for suspicious,
blocked, and error outcomes. This API property narrows common integration
mistakes, but it cannot prevent a separate application path from retaining or
opening the original input. The embedding application and its filesystem
permissions remain part of the trusted computing base.
